\documentclass[letterpaper,12pt]{article}
\usepackage[margin=1in]{geometry}
\usepackage{graphicx}
\usepackage{epstopdf}
\usepackage{hyperref}
%\hypersetup{colorlinks=false}
\hypersetup{colorlinks,citecolor=black,filecolor=red,linkcolor=black,urlcolor=black,pdfnewwindow=true,bookmarksnumbered}
\usepackage{amsmath}
\usepackage{amsxtra}
\usepackage{amstext}
\usepackage{amssymb}
\usepackage{latexsym}
% The `comment' package is not installed in this TeX distribution; the
% `verbatim' package (which is) provides a robust comment environment that
% swallows its body, including verbatim-like content.
\usepackage{verbatim}

\title{\vspace{-0.5 in} APPH 6102\\ Project \#1 \vspace{-0.5 in}}
\date{distributed: Feb 2026 \\ credit: 15 points}
%\author{C. Paz-Soldan}

\begin{document}

\maketitle
\vspace{-25 pt}
In these project components you will be provided with MATLAB and python scripts to get started. The deliverables are: 1) a copy (zipped) of your scripts that can be run without additional files. 2) a slide show presentation of your results, including figures and short descriptive text (like a research presentation. 3) a zoom recording of you delivering your slide show, target 10-15 minutes (short and sweet).\\

\underline{Part \#1: Proto-Cleo Stellarator}\\

In this part you will numerically trace out magnetic surfaces in a full 3D geometry and derive some properties of the field. The field structure is the “Proto-Cleo” stellarator, a 3-fold poloidal and 7-fold toroidal periodicity device. It operated at the UK Culham Laboratory c. 1970\footnote{A summary of the device can be found here: \url{https://doi.org/10.1063/1.1693643}.}. The Proto-Cleo device had a major radius of 0.4 m and a plasma radius of ~0.05 m.\\

MATLAB/python scripts have been provided to calculate the helical field from the Biot-Savart Law from the original coil geometry. Note there is also a toroidal field added analytically (the full system had helical coils and toroidal field coils). You have also been provided with a field line integrator for actually stepping along the magnetic field lines. The assignment has several components, enumerated below.\\

1) By line tracing find the magnetic axis (there may be more than one!), find the outermost separatrix where well-defined surfaces turn into divergent field lines, and plot at least eight magnetic surfaces between the axis and the separatrix. Plot these at 0, 1/4, 1/2 and 3/4 field period toroidal cuts. Hint: Plot Poincare sections.\\

2) Caulcuate \& plot the rotational transform vs minor radius. Hint: Count transits.\\

3) Plot $\left|B\right|$ vs length $L$ for a field line at roughly 3/4 of the way to the separatrix from the axis. Follow this for a sufficient distance that you can see both the toroidal and helical modulation of the magnetic field strength.\\

4) Add a 3/1 `error' field of the form $b_r=b_0\cos{(\phi{})}\cos{(3\theta)}$ to the nominal magnetic field. This error should resonant with the q=3 surface, and islands should appear. Try $b_0=0.5$ Gauss and plot the islands. Vary $b_0$ over a range of values. Can you determine how the island width scales? When are the surfaces completely destroyed? What happens away from the islands?

% Alternative to the Solovev equilibrium project. Let's see how hard this is.
\begin{comment}
\newpage
\clearpage

\underline{Part \#2: Tokamak}\\
Lets follow some orbits in a tokamak geometry. You are provided with a MATLAB routine that will calculate the magnetic field vector in a large aspect ratio tokamak with q ~3 at the R=10 m radius. You are also provided with routines that will calculate the gradient of the magnetic field strength along a field line, and in 3D. Look at the field-line tracer in Part \# and modify it for this problem. It should use MATLAB's `ode45' routine.\\

1) Trace out orbits for co and counter injected circulating ions with energies of 500 eV. \\

2) Calculate the necessary pitch angle for a particle which will turn around at R = 10 m when launched at the R = 10.3 m Z = 0 m point, and provide a trace of your simulation of that orbit.
\end{comment}

\newpage
\clearpage

\underline{Part \#2: Toroidal Equilibrium Magnetic Sensitivity}\\

This project will use existing Solovev solutions to the Grad-Shafranov equation to see how one might place magnetic measurements around the tokamak to reconstruct the equilibrium. The project requires you to download the `One Size Fits All' Grad-Shafranov MATLAB solver by A. Cerfon et al of NYU\footnote{The solver used to be online, but now it's only on courseworks}. We will use the limited plasma tool. Your progress will be facilitated by converting the `main.m' MATLAB solver into a function that can be placed inside a loop. \\

0) Read the `One Size Fits All' paper to understand how the solutions are generated. Note its basic similarity to Homework \#2.\\

1) Use Cerfon's tool to define a base limited (no X-point) configuration. Options include: TFTR, JET, ITER, NSTX, MAST, FRC, Spheromak.\\

2) Define a rectangular surface (height $h$, width $w$) that is outside the plasma, encircling it but not intersecting it. Hint: avoid R=0 singularities, and you may need to expand the $\Psi$ computational domain. Pretend this is the vacuum vessel. \\

3) Use your knowledge of toroidal equilibrium to numerically convert the output $\Psi$ into $B_R$ and $B_Z$ components.\\

4) Find and plot the flux and field components along your rectangular surface. It may be easier to treat each side separately.\\

5) For the configuration you chose, repeat \#4 the above but with large reductions in elongation (make sure it still fits in the surface of part \#2.\\

6) Based on the variations found in \#5, where would you position $B_R$, $B_Z$, and $\Psi$ sensors to maximize their sensitivity to the plasma shape change you imposed? Hint: Look where the signal is large and where it is changing the most.\\

7) Repeat the above but varying the trianglarity (delta) from zero to above the reference shape. Where would $B_R$, $B_Z$, and $\Psi$ sensors be ideally located ?\\

8) Repeat the above but varying the `A' parameter from zero to one in the reference shape. Where would $B_R$, $B_Z$, and $\Psi$ sensors be ideally located ?

\end{document}
